\documentclass{article}
\usepackage{amsmath}
\usepackage{amsfonts}
\usepackage{amscd}
\usepackage{amssymb}
\usepackage{xcolor}
\usepackage{listings}

\lstdefinestyle{colored}{
    backgroundcolor=\color{lightgray}, % color de fondo
    frame=single,                     % marco simple
    rulecolor=\color{lime},          % color del marco
    basicstyle=\ttfamily\small,      % fuente del código
}

\title{\Huge Modelos para Variables Aleatorias}
\author{}
\date{}

\begin{document}
\maketitle

\section*{\Large Modelos para Variables Aleatorias Discretas}

1. \textbf{\Large{Distribución de Bernoulli}}:
   \begin{itemize}
       \item \textbf{Descripción}: Representa un experimento con dos resultados posibles (éxito o fracaso).
       \item \textbf{Parámetro}: \( p \) (probabilidad de éxito).
       \item \textbf{Función de Masa de Probabilidad (PMF)}:
       \[
       P(X=1) = p, \quad P(X=0) = 1 - p
       \]
       \item \textbf{Valor Esperado}:
       \[
       E[X] = p
       \]
       \item \textbf{Media}:
       \[
       \mu_X = p
       \]
        \item \textbf{Varianza}:
       \[
       \sigma = p(1-p)
       \]
       \item \textbf{Desviación Estándar}:
       \[
       \sigma = \sqrt{p(1-p)}
       \]
       \item \textbf{Función de Distribución Acumulativa (CDF)}:
       \[
       F(x) = 
       \begin{cases} 
       0 & \text{si } x < 0 \\ 
       1 - p & \text{si } 0 \leq x < 1 \\ 
       1 & \text{si } x \geq 1 
       \end{cases}
       \]
       \item \textbf{R Commands}:
      
           \begin{lstlisting}[style=colored, language=R]
           # PMF (funcion de masa de probabilidad)
           dbinom(x, size=1, prob=p)
           
           # CDF (funcion de distribucion acumulada)
           pbinom(x, size=1, prob=p)
           
           # Para generar simulacion
           rbinom(n, size=1, prob=p)       
          \end{lstlisting}
       
   \end{itemize}

2. \textbf{{\Large Distribución Binomial}}:
   \begin{itemize}
       \item \textbf{Descripción}: Modelo de \( n \) ensayos independientes de Bernoulli.
       \item \textbf{Parámetros}: \( n \) (número de ensayos), \( p \) (probabilidad de éxito).
       \item \textbf{Función de Masa de Probabilidad (PMF)}:
       \[
       P(X=k) = \binom{n}{k} p^k (1-p)^{n-k}, \quad k = 0, 1, \ldots, n
       \]
       \item \textbf{Valor Esperado}:
       \[
       E[X] = n \cdot p
       \]
       \item \textbf{Media}:
       \[
       \text{Media} = n \cdot p
       \]
       \item \textbf{Desviación Estándar}:
       \[
       \sigma = \sqrt{n \cdot p \cdot (1-p)}
       \]
       \item \textbf{Función de Distribución Acumulativa (CDF)}:
       \[
       F(x) = P(X \leq k) = \sum_{i=0}^{\lfloor x \rfloor} \binom{n}{i} p^i (1-p)^{n-i}
       \]
       \item \textbf{R Commands}:
        \begin{lstlisting}[style=colored, language=R]
       # PMF
       dbinom(k, size=n, prob=p)
       
       # CDF
       pbinom(k, size=n, prob=p)
       
       # Simulation
       rbinom(n, size=n, prob=p)
       \end{lstlisting}
   \end{itemize}

3. \textbf{{\Large Distribución de Poisson}}:
   \begin{itemize}
       \item \textbf{Descripción}: Modelo de eventos que ocurren en un intervalo fijo de tiempo o espacio.
       \item \textbf{Parámetro}: \( \lambda \) (tasa promedio de ocurrencias).
       \item \textbf{Función de Masa de Probabilidad (PMF)}:
       \[
       P(X=k) = \frac{\lambda^k e^{-\lambda}}{k!}, \quad k = 0, 1, 2, \ldots
       \]
       \item \textbf{Valor Esperado}:
       \[
       E[X] = \lambda
       \]
       \item \textbf{Media}:
       \[
       \text{Media} = \lambda
       \]
       \item \textbf{Desviación Estándar}:
       \[
       \sigma = \sqrt{\lambda}
       \]
       \item \textbf{Función de Distribución Acumulativa (CDF)}:
       \[
       F(x) = P(X \leq k) = \sum_{i=0}^{\lfloor x \rfloor} \frac{\lambda^i e^{-\lambda}}{i!}
       \]
       \item \textbf{R Commands}:
       \begin{lstlisting}[style=colored, language=R]
       # PMF
       dpois(x, lambda)
       
       # CDF
       ppois(x, lambda)
       
       # Simulation
       rpois(n, lambda)
       \end{lstlisting}
   \end{itemize}

4. \textbf{{\Large Distribución Geométrica}}:
   \begin{itemize}
       \item \textbf{Descripción}: Modelo de la cantidad de ensayos necesarios para obtener el primer éxito.
       \item \textbf{Parámetro}: \( p \) (probabilidad de éxito en cada ensayo).
       \item \textbf{Función de Masa de Probabilidad (PMF)}:
       \[
       P(X=k) = (1-p)^{k-1} p, \quad k = 1, 2, \ldots
       \]
       \item \textbf{Valor Esperado}:
       \[
       E[X] = \frac{1}{p}
       \]
       \item \textbf{Media}:
       \[
       \text{Media} = \frac{1}{p}
       \]
       \item \textbf{Desviación Estándar}:
       \[
       \sigma = \frac{\sqrt{1-p}}{p}
       \]
       \item \textbf{Función de Distribución Acumulativa (CDF)}:
       \[
       F(x) = P(X \leq k) = 1 - (1-p)^k, \quad k = 1, 2, \ldots
       \]
       \item \textbf{R Commands}:
       \begin{lstlisting}[style=colored, language=R]
       # PMF
       dgeom(k-1, prob=p)
       
       # CDF
       pgeom(k-1, prob=p)
       
       # Simulation
       rgeom(n, prob=p)
       \end{lstlisting}
   \end{itemize}

5. \textbf{{ \Large Distribución Hipergeométrica}}:
   \begin{itemize}
       \item \textbf{Descripción}: Modelo para muestreo sin reemplazo de una población finita.
       \item \textbf{Parámetros}: \( N \) (tamaño de la población), \( K \) (número de éxitos en la población), \( n \) (tamaño de la muestra).
       \item \textbf{Función de Masa de Probabilidad (PMF)}:
       \[
       P(X=k) = \frac{\binom{K}{k} \binom{N-K}{n-k}}{\binom{N}{n}}, \quad k = 0, 1, \ldots, \min(K, n)
       \]
       \item \textbf{Valor Esperado}:
       \[
       E[X] = \frac{nK}{N}
       \]
       \item \textbf{Media}:
       \[
       \text{Media} = \frac{nK}{N}
       \]
       \item \textbf{Desviación Estándar}:
       \[
       \sigma = \sqrt{\frac{nK(N-K)(N-n)}{N^2(N-1)}}
       \]
       \item \textbf{Función de Distribución Acumulativa (CDF)}:
       \[
       F(x) = P(X \leq k) = \sum_{j=0}^{\lfloor x \rfloor} \frac{\binom{K}{j} \binom{N-K}{n-j}}{\binom{N}{n}}
       \]
       \item \textbf{R Commands}:
       \begin{lstlisting}[style=colored, language=R]
       # PMF
       dhyper(k, m=K, n=N-K, size=n)
       
       # CDF
       phyper(k, m=K, n=N-K, size=n)
       
       # Simulation
       rhyper(n, m=K, n=N-K, size=n)
       \end{lstlisting}
   \end{itemize}

6. \textbf{{ \Large Distribución Negativa Binomial}}:
   \begin{itemize}
      

 \item \textbf{Descripción}: Modelo de la cantidad de ensayos necesarios para obtener un número fijo de éxitos (generaliza la distribución geométrica).
       \item \textbf{Parámetros}: \( r \) (número de éxitos deseados), \( p \) (probabilidad de éxito).
       \item \textbf{Función de Masa de Probabilidad (PMF)}:
       \[
       P(X=k) = \binom{k+r-1}{r-1} p^r (1-p)^k, \quad k = 0, 1, 2, \ldots
       \]
       \item \textbf{Valor Esperado}:
       \[
       E[X] = \frac{r(1-p)}{p}
       \]
       \item \textbf{Media}:
       \[
       \text{Media} = \frac{r(1-p)}{p}
       \]
       \item \textbf{Desviación Estándar}:
       \[
       \sigma = \sqrt{\frac{r(1-p)}{p^2}}
       \]
       \item \textbf{Función de Distribución Acumulativa (CDF)}:
       \[
       F(x) = P(X \leq k) = \sum_{j=0}^{\lfloor x \rfloor} \binom{j+r-1}{r-1} p^r (1-p)^j
       \]
       \item \textbf{R Commands}:
       \begin{lstlisting}[style=colored, language=R]
       # PMF
       dnbinom(k, size=r, prob=p)
       
       # CDF
       pnbinom(k, size=r, prob=p)
       
       # Simulacion
       rnbinom(n, size=r, prob=p)
       \end{lstlisting}
   \end{itemize}

7. \textbf{{ \Large Distribución Multinomial}}:
   \begin{itemize}
       \item \textbf{Descripción}: Extensión de la distribución binomial a más de dos resultados.
       \item \textbf{Parámetros}: \( n \) (número de ensayos), \( p_1, p_2, \ldots, p_k \) (probabilidades de cada resultado).
       \item \textbf{Función de Masa de Probabilidad (PMF)}:
       \[
       P(X_1=k_1, X_2=k_2, \ldots, X_k=k_k) = \frac{n!}{k_1! k_2! \ldots k_k!} p_1^{k_1} p_2^{k_2} \ldots p_k^{k_k}
       \]
       \item \textbf{Valor Esperado}:
       \[
       E[X_i] = n \cdot p_i, \quad i = 1, 2, \ldots, k
       \]
       \item \textbf{Media}:
       \[
       \text{Media} = n \cdot p_i \text{ para cada } i
       \]
       \item \textbf{Desviación Estándar}:
       \[
       \sigma_i = \sqrt{n \cdot p_i \cdot (1 - p_i)}, \quad i = 1, 2, \ldots, k
       \]
       \item \textbf{Función de Distribución Acumulativa (CDF)}: (No existe una fórmula general simple, se utiliza el enfoque de simulación)
       \item \textbf{R Commands}:
       \begin{lstlisting}[style=colored, language=R]
       # PMF
       dmultinom(c(k1, k2, ..., kk), size=n, prob=c(p1, p2, ..., pk))
       
       # CDF (no hay funcion directa)
       
       # Simulacion
       rmultinom(n, size=n, prob=c(p1, p2, ..., pk))
       \end{lstlisting}
   \end{itemize}

\end{document}
